\documentclass[12pt,a4paper]{article}
\usepackage[utf8]{inputenc}
\usepackage{graphicx}
\usepackage{amsmath}
\usepackage{hyperref}
\usepackage{booktabs}
\usepackage{listings}
\usepackage{xcolor}
\usepackage{geometry}
\geometry{margin=1in}

\title{\textbf{Solar Power Generation Forecasting Using Machine Learning} \\ \large A Time-Series Prediction System}
\author{Your Name \\ Roll Number \\ AI/ML Course Project}
\date{\today}

\begin{document}

\maketitle
\newpage

\tableofcontents
\newpage

\section{Abstract}
Solar energy forecasting is essential for efficient grid management and renewable energy integration. This project develops a machine learning-based forecasting system to predict solar power generation using the Anikannal Solar Power Generation Dataset. Two models were implemented and compared: Linear Regression and Random Forest Regressor. The system achieves a Mean Absolute Error (MAE) of 17.26 kW using Random Forest, demonstrating effective prediction capability. A web-based interface was developed using Streamlit for real-time predictions.

\section{Introduction}

\subsection{Background}
Solar power generation is inherently variable, depending on weather conditions, time of day, and seasonal patterns. Accurate forecasting enables:
\begin{itemize}
    \item Optimal grid load balancing
    \item Efficient energy storage management
    \item Reduced reliance on backup power sources
    \item Better integration of renewable energy
\end{itemize}

\subsection{Problem Statement}
Given historical weather data and power generation records, predict the AC power output of a solar plant at 15-minute intervals.

\subsection{Objectives}
\begin{enumerate}
    \item Develop time-series forecasting models for solar power prediction
    \item Compare Linear Regression and Random Forest performance
    \item Engineer relevant temporal and weather-based features
    \item Deploy an interactive web application for real-time predictions
\end{enumerate}

\section{Literature Review}
Time-series forecasting for solar energy has been extensively studied using various approaches:
\begin{itemize}
    \item \textbf{Statistical Methods}: ARIMA, SARIMA for capturing temporal patterns
    \item \textbf{Machine Learning}: Random Forest, Gradient Boosting for non-linear relationships
    \item \textbf{Deep Learning}: LSTM, GRU networks for sequential dependencies
\end{itemize}

This project focuses on traditional machine learning approaches due to their interpretability and computational efficiency.

\section{Dataset Description}

\subsection{Data Sources}
The Anikannal Solar Power Generation Dataset (Plant 1) consists of two primary files:

\subsubsection{Generation Data}
\begin{itemize}
    \item \textbf{Frequency}: 15-minute intervals
    \item \textbf{Records}: 34,000+ observations
    \item \textbf{Key Features}: DATE\_TIME, AC\_POWER (target), DC\_POWER, DAILY\_YIELD
\end{itemize}

\subsubsection{Weather Data}
\begin{itemize}
    \item \textbf{Frequency}: Hourly measurements
    \item \textbf{Records}: 3,000+ observations
    \item \textbf{Key Features}: AMBIENT\_TEMPERATURE, MODULE\_TEMPERATURE, IRRADIATION
\end{itemize}

\subsection{Data Characteristics}
\begin{table}[h]
\centering
\begin{tabular}{@{}lcc@{}}
\toprule
\textbf{Feature} & \textbf{Range} & \textbf{Unit} \\ \midrule
AC\_POWER & 0 - 1410.95 & kW \\
AMBIENT\_TEMPERATURE & 15 - 45 & °C \\
MODULE\_TEMPERATURE & 10 - 65 & °C \\
IRRADIATION & 0 - 1.22 & kW/m² \\ \bottomrule
\end{tabular}
\caption{Dataset Feature Ranges}
\label{tab:features}
\end{table}

\section{Methodology}

\subsection{Data Preprocessing}

\subsubsection{Time-Aware Merging}
Since generation data is recorded every 15 minutes and weather data hourly, a time-aware merge strategy was employed:

\begin{lstlisting}[language=Python]
merged_df = pd.merge_asof(gen_df, weather_df, 
                          on='DATE_TIME', 
                          direction='nearest')
\end{lstlisting}

This approach matches each generation record with the nearest weather measurement, preserving temporal relationships.

\subsubsection{Data Cleaning}
\begin{enumerate}
    \item \textbf{Night Period Removal}: Filtered out records where AC\_POWER = 0 (non-generation periods)
    \item \textbf{Duplicate Removal}: Eliminated redundant timestamps
    \item \textbf{Missing Value Handling}: Dropped rows with NaN values after feature engineering
\end{enumerate}

\subsection{Feature Engineering}

\subsubsection{Temporal Features}
Extracted time-based features to capture daily and seasonal patterns:
\begin{align}
\text{HOUR} &= \text{hour}(\text{DATE\_TIME}) \quad \in [0, 23] \\
\text{MONTH} &= \text{month}(\text{DATE\_TIME}) \quad \in [1, 12] \\
\text{DAY\_OF\_WEEK} &= \text{dayofweek}(\text{DATE\_TIME}) \quad \in [0, 6]
\end{align}

\subsubsection{Time-Series Features}
Introduced lag and rolling features to capture temporal dependencies:

\textbf{Lag Feature:}
\begin{equation}
\text{lag\_1}_t = \text{AC\_POWER}_{t-1}
\end{equation}

\textbf{Rolling Mean:}
\begin{equation}
\text{rolling\_mean\_3}_t = \frac{1}{3}\sum_{i=0}^{2} \text{AC\_POWER}_{t-i}
\end{equation}

\subsubsection{Final Feature Set}
\begin{equation}
X = [T_{amb}, T_{mod}, I, H, M, D, L_1, R_3]
\end{equation}

Where:
\begin{itemize}
    \item $T_{amb}$ = Ambient Temperature
    \item $T_{mod}$ = Module Temperature
    \item $I$ = Irradiation
    \item $H$ = Hour
    \item $M$ = Month
    \item $D$ = Day of Week
    \item $L_1$ = Lag-1 feature
    \item $R_3$ = Rolling mean (window=3)
\end{itemize}

\subsection{Train-Test Split}
A chronological 80-20 split was performed without shuffling to maintain time-series integrity:

\begin{equation}
\text{Split Index} = \lfloor 0.8 \times N \rfloor
\end{equation}

Where $N$ is the total number of samples. This ensures the model is trained on past data and tested on future data, simulating real-world forecasting scenarios.

\subsection{Feature Scaling}
StandardScaler was applied to features for Linear Regression:

\begin{equation}
X_{scaled} = \frac{X - \mu}{\sigma}
\end{equation}

Where $\mu$ is the mean and $\sigma$ is the standard deviation. Random Forest does not require scaling as it uses tree-based splitting.

\subsection{Model Development}

\subsubsection{Linear Regression}
A baseline linear model was implemented:

\begin{equation}
\hat{y} = \beta_0 + \sum_{i=1}^{8} \beta_i x_i
\end{equation}

Where $\hat{y}$ is the predicted AC power, $\beta_0$ is the intercept, and $\beta_i$ are the coefficients for each feature.

\subsubsection{Random Forest Regressor}
An ensemble model with 100 decision trees:

\begin{equation}
\hat{y} = \frac{1}{N_{trees}}\sum_{i=1}^{N_{trees}} f_i(X)
\end{equation}

Where $f_i(X)$ is the prediction from the $i$-th tree.

\textbf{Hyperparameters:}
\begin{itemize}
    \item n\_estimators = 100
    \item random\_state = 42
    \item n\_jobs = -1 (parallel processing)
\end{itemize}

\section{Evaluation Metrics}

\subsection{Mean Absolute Error (MAE)}
\begin{equation}
\text{MAE} = \frac{1}{n}\sum_{i=1}^{n}|y_i - \hat{y}_i|
\end{equation}

MAE represents the average absolute difference between predicted and actual values, providing an intuitive measure of prediction accuracy in kW.

\subsection{Root Mean Squared Error (RMSE)}
\begin{equation}
\text{RMSE} = \sqrt{\frac{1}{n}\sum_{i=1}^{n}(y_i - \hat{y}_i)^2}
\end{equation}

RMSE penalizes larger errors more heavily, making it sensitive to outliers.

\section{Results and Analysis}

\subsection{Model Performance}

\begin{table}[h]
\centering
\begin{tabular}{@{}lcc@{}}
\toprule
\textbf{Model} & \textbf{MAE (kW)} & \textbf{RMSE (kW)} \\ \midrule
Linear Regression & 17.33 & 32.09 \\
Random Forest & \textbf{17.26} & \textbf{32.45} \\ \bottomrule
\end{tabular}
\caption{Model Performance Comparison}
\label{tab:results}
\end{table}

\subsection{Key Findings}
\begin{enumerate}
    \item \textbf{Random Forest Superiority}: Achieved slightly lower MAE (17.26 kW vs 17.33 kW)
    \item \textbf{Low Error Rate}: Average error of ~17 kW on a range of 0-1410 kW represents approximately 1-5\% error
    \item \textbf{Feature Importance}: Irradiation emerged as the most influential feature, followed by module temperature and lag features
    \item \textbf{Temporal Patterns}: Hour and month features effectively captured daily and seasonal variations
\end{enumerate}

\subsection{Model Comparison Analysis}
\begin{itemize}
    \item \textbf{Linear Regression}: Faster training, interpretable coefficients, but assumes linear relationships
    \item \textbf{Random Forest}: Captures non-linear patterns (e.g., efficiency drop at extreme temperatures), more robust to outliers
\end{itemize}

\section{Deployment}

\subsection{Web Application}
A Streamlit-based web interface was developed with the following features:

\begin{itemize}
    \item \textbf{Model Selection}: Toggle between Linear Regression and Random Forest
    \item \textbf{Interactive Inputs}: Sliders and number inputs for all 8 features
    \item \textbf{Real-time Prediction}: Instant power output prediction
    \item \textbf{Visualizations}: 
    \begin{itemize}
        \item Gauge chart for predicted power
        \item Model comparison metrics
        \item Feature value bar chart
    \end{itemize}
    \item \textbf{Sample Scenarios}: Pre-configured examples (peak, morning, evening)
\end{itemize}

\subsection{Deployment Architecture}
\begin{verbatim}
User Interface (Streamlit)
    |
    v
Model Loading (pickle)
    |
    v
Feature Scaling (StandardScaler)
    |
    v
Prediction (Linear Regression / Random Forest)
    |
    v
Visualization (Plotly)
\end{verbatim}

\section{Challenges and Solutions}

\begin{table}[h]
\centering
\begin{tabular}{@{}p{5cm}p{7cm}@{}}
\toprule
\textbf{Challenge} & \textbf{Solution} \\ \midrule
Mismatched time intervals (15-min vs hourly) & Used merge\_asof for time-aware merging \\
Night periods with zero power & Filtered AC\_POWER = 0 records \\
Feature scale differences & Applied StandardScaler for Linear Regression \\
Temporal dependencies & Introduced lag and rolling mean features \\
Data leakage risk & Chronological split without shuffling \\ \bottomrule
\end{tabular}
\caption{Challenges and Solutions}
\label{tab:challenges}
\end{table}

\section{Conclusion}

This project successfully developed a machine learning-based solar power forecasting system with the following achievements:

\begin{enumerate}
    \item \textbf{Accurate Predictions}: Achieved MAE of 17.26 kW using Random Forest
    \item \textbf{Effective Feature Engineering}: Temporal and time-series features significantly improved model performance
    \item \textbf{Model Comparison}: Demonstrated trade-offs between Linear Regression and Random Forest
    \item \textbf{Practical Deployment}: Developed user-friendly web interface for real-time predictions
\end{enumerate}

\subsection{Real-World Impact}
\begin{itemize}
    \item Enables grid operators to anticipate solar generation fluctuations
    \item Facilitates optimal energy storage scheduling
    \item Reduces dependency on fossil fuel backup systems
    \item Supports renewable energy integration goals
\end{itemize}

\section{Future Work}

\subsection{Model Enhancements}
\begin{itemize}
    \item \textbf{Deep Learning}: Implement LSTM/GRU networks for better sequential pattern recognition
    \item \textbf{Ensemble Methods}: Combine multiple models using stacking or blending
    \item \textbf{Hyperparameter Optimization}: Grid search or Bayesian optimization for Random Forest
\end{itemize}

\subsection{Feature Additions}
\begin{itemize}
    \item Weather forecasts (cloud cover, humidity, wind speed)
    \item Historical weather patterns
    \item Equipment degradation factors
    \item Seasonal decomposition features
\end{itemize}

\subsection{Deployment Improvements}
\begin{itemize}
    \item Multi-step ahead forecasting (predict next 24 hours)
    \item Confidence intervals for predictions
    \item Automated model retraining pipeline
    \item Integration with real-time weather APIs
\end{itemize}

\section{References}

\begin{enumerate}
    \item Anikannal Solar Power Generation Dataset, Kaggle
    \item Pedregosa et al. (2011). Scikit-learn: Machine Learning in Python. JMLR 12, pp. 2825-2830
    \item Breiman, L. (2001). Random Forests. Machine Learning 45(1), pp. 5-32
    \item Streamlit Documentation: https://docs.streamlit.io
    \item Pandas Documentation: https://pandas.pydata.org
\end{enumerate}

\section{Appendix}

\subsection{Code Repository}
Project code available at: [GitHub URL]

\subsection{Hosted Application}
Live demo: [Streamlit Cloud URL]

\subsection{System Requirements}
\begin{itemize}
    \item Python 3.8+
    \item Libraries: streamlit, pandas, numpy, scikit-learn, plotly
    \item RAM: 2GB minimum (for Random Forest model)
    \item Storage: 300MB (including dataset and models)
\end{itemize}

\end{document}
